\newif\ifuseseminar
\useseminartrue
\input{../../latex_common/header.tex.frag}

\title{Relatividade Restrita -- Prova: Funções de Green e Álgebra Geométrica}

\newcommand{\e}{\mathsf{e}}

\begin{document}

\begin{enumerate}

      \item Bivetores e Boosts

            \begin{enumerate}
                  \item Calcule o quadrado dos seguintes bivetores e classifique cada um
                        como tipo-tempo, tipo-espaço ou nulo:
                        $$(\e^0 \wedge \e^1)^2, \quad (\e^1 \wedge \e^2)^2, \quad (\e^0 \wedge \e^2)^2.$$

                  \item Considere o bivetor $B = \e^0 \wedge \e^1$ e o rotor $R =
                              e^{\eta B/2}$, com parâmetro de rapidez $\eta$. Mostre que a
                        transformação
                        $$u \mapsto R u R^\mathrm{T},$$ onde $R^\mathrm{T} = e^{-\eta
                                          B/2}$, corresponde a um boost de Lorentz ao longo da
                        direção $x^1$ para um vetor qualquer $u = u^\mu \e_\mu$.
                        Interprete o resultado.
                  \item Mostre que se $B = \e^1 \wedge \e^2$ e $R = e^{\theta B/2}$, então
                        $$u \mapsto R u R^\mathrm{T}$$ corresponde a uma rotação de
                        $\theta$ radianos em torno do eixo $\e^3$.
            \end{enumerate}

      \item Pseudoscalar e Dualidade

            \begin{enumerate}
                  \item Mostre que o pseudoscalar $I=\e^0 \wedge \e^1 \wedge \e^2 \wedge
                              \e^3$ anti-comuta com cada vetor base, ou seja,
                        $$I \e^\mu = - \e^\mu I, \quad \mu = 0,1,2,3.$$

                  \item Usando o pseudoscalar, calcule o dual do bivetor $B = \e^1
                              \wedge \e^2$ ou seja $IB$ e interprete o significado
                        geométrico desse dual.
            \end{enumerate}

      \item Projeções e Componentes Temporais

            \begin{enumerate}
                  \item Mostre que a componente temporal de $x = x^\mu \e_\mu$ em
                        relação a $w = \e_0$ pode ser recuperada por
                        $$ct = - x \cdot w.$$

                  \item Defina a projeção espacial ortogonal a $w$ por
                        $$x_\perp = \frac{x + w x w}{2}$$
                        e prove que
                        $$x_\perp \cdot w = 0,\qquad\text{e}\quad wx_\perp = -x_\perp w.$$
            \end{enumerate}

      \item Equação da Onda e Função de Green

            \begin{enumerate}
                  \item Defina a função de Green $G(t, \vec{x}; t', \vec{x}')$ associada
                        à equação da onda
                        $$\Box \phi(t, \vec{x}) = \rho(t, \vec{x}),$$
                        onde
                        $$\Box = -\frac{1}{c^2} \frac{\partial^2}{\partial t^2} +
                              \nabla^2,$$ e explique sua interpretação física, destacando as
                        condições de causalidade para as funções de Green retardada e
                        avançada.

                  \item A função de Green retardada em 3 dimensões espaciais pode ser
                        expressa como
                        $$
                              G_{\mathrm{ret}}(ct - ct', \vec{x} - \vec{x}') =
                              \frac{\delta\left(ct - ct' - |\vec{x} - \vec{x}'|\right)}{4\pi |\vec{x} - \vec{x}'|}.
                        $$
                        Explique o significado físico do suporte desta função estar
                        limitado à superfície do cone de luz.

                  \item Calcule
                        $$
                              \phi(t,\vec{x}) = \int G_{\mathrm{ret}}(t, \vec{x}; t', \vec{x}') \, \delta^3(\vec{x} - \vec{x}') \, \dd t' \dd^3 \vec{x}',
                        $$
                        e explique o significado físico desse resultado.
            \end{enumerate}

\end{enumerate}

\end{document}
